\chapter{Grundlagen}
\label{ch:grundlagen}

\section{Was Hefe eigentlich macht}
Während der Gärung wandelt die Hefe \emph{Saccharomyces cerevisiae} vergärbaren
Zucker in Ethanol und Kohlendioxid um:
\begin{equation}
  \mathrm{C_6H_{12}O_6} \;\longrightarrow\; 2\,\mathrm{C_2H_5OH} + 2\,\mathrm{CO_2}
  \label{eq:gaerung}
\end{equation}
Nebenbei entstehen hunderte Aromastoffe. Besonders wichtig sind:
\begin{itemize}
  \item \textbf{Ester} wie Isoamylacetat (Banane) und Ethylhexanoat (Apfel),
  \item \textbf{höhere Alkohole} wie Isoamylalkohol (lösungsmittelartig),
  \item \textbf{Phenole} wie 4-Vinylguajakol (Gewürznelke, typisch für Weizenbier).
\end{itemize}

\section{Temperatur als Stellschraube}
Die Esterbildung folgt näherungsweise einer Arrhenius-Beziehung
\citep{gerstner2020kinetik}:
\begin{equation}
  k(T) = A \cdot \exp\!\left(-\frac{E_a}{R\,T}\right)
  \label{eq:arrhenius}
\end{equation}
Nach \cref{eq:arrhenius} sollte eine höhere Temperatur die Bildung deutlich
beschleunigen -- wie stark, ist für obergärige Hefen jedoch kaum dokumentiert
\citep{malz2021ester, hefeling2018obergaerig}.
